% 第 5 章 结论
% 本章引用键（实际使用）：fips203, kyber2018, kyberslash, clangover, dudect, rtlf, binsec
% 本章不放图、不放表；数字回指第 3 章与附录 D
% 数字来源：5.1 用 A-01–A-09、A-15、A-17、A-18、A-22–A-24、A-27、A-28；5.2 用 C-01–C-15 与 A-29、A-30；5.3 为计划，不含实测数字

\section{结论}

\subsection{主要结论}

本作品在 ML-KEM 解封装上建立了一条“检测—定位—防护—复测”的分层证据链，并在真实历史案例、人工校准与阴性对照上把四段逐一走通。五条判断各自绑定证据层与适用边界；数字回指第 3 章与附录 D。% A-01 % A-02 % A-03 % A-22 % A-24

最直接的一条判断是：数十纳秒量级的真实泄漏可以被检出，也可以在产物层核验修复。M1 未修版本在主类对上的合并均值差为 $+32.311$ ns，95\% 区间 [20.282, 44.339]，预注册 Holm 家族中 $p=1.30\times10^{-5}$，低于阈值 0.00833；族外的 dudect\cite{dudect} 与 RTLF\cite{rtlf} 给出同向结论。官方修复后合并差为 $-8.150$ ns，区间 [$-21.916$, 5.616]，三次独立运行的区间都含 0；解封装开销增加 1.51\%。这条判断限定在本机：32.311 ns 只有 2--3 个分辨率单位，依赖每类 10,000 次、三次运行与预注册家族，换产物或换机器要重新判定。% A-01 % A-02 % A-09 % A-11 % A-25 % A-30

第二条结论：定位层与计时层必须分开登记。M2 的三处秘密除法在定位层稳定出现，可变时延检查报 1,280 errors / 3 contexts，官方修复后归零；同一产物的端到端合并差为 $+4.58$ ns，区间 [$-3.08$, 12.25]，除法热点微基准为 $+0.044$ ns，区间 [$-0.430$, 0.517]，两层都未检出。M2 的检测预算为 25--67 ns，大于除法可能产生的效应；-O2 与 -O3 的产物里秘密分支的命中位置更多，计时同样未检出。定位阳性回答“机制在不在”，计时阳性回答“本机分不分得开”，两者不能互相代替。% A-03 % A-04 % A-15 % A-18 % A-26

修复已知泄漏不等于产物达到常数时间。90 格排查的结果是 5 格存在、75 格不存在、10 格不适用、0 格未检查，覆盖三个实现、五个解封装步骤与六类泄漏来源。KyberSlash\cite{kyberslash} 官方修复提交只消除了除法，消息编码函数里与 clangover\cite{clangover} 同源的分支仍在；未修与已修两侧的指令字节逐字节相同，秘密标记各报 256 errors / 1 context。这是历史产物的机制层观察，未做逐指令时序归因；结论停在机制层，不属于新漏洞。同一份源码的常数时间承诺，只有落到具体产物、具体证据层与具体检测预算上才有内容；修复一条机制，不自动清除同一条路径上的其他机制。% A-22 % A-23

阴性对照没有出现误报。M5 在本机、本编译产物、本类对与每类 10,000 次的预算下，合并差为 $-5.30$ ns，区间 [$-11.77$, 1.16] 含 0，预注册家族没有被拒绝；在 M5 上运行过的 6 种工具误报为 0。阴性对照检验的是工具与流程会不会把无泄漏读成泄漏，不是给 M5 发常数时间证书；这条阴性只排除大于该预算的效应，不排除更小的效应，也不覆盖其他产物。分母是 6 个工具与 1 个对象，结论只覆盖本次矩阵。% A-08 % A-24

测量条件决定可检出性。锁频的价值在确定性，不在降噪：不锁频与锁频的变异系数比为 0.964、1.006、1.089，方向不一致，锁频提供的是可核对、可恢复的工作点。M1 的泄漏机制在 aarch64 上不复现，同一份源码被编成位域提取与条件选择；同机网络五档在本预算下全部未检出，分辨同一量级效应所需的样本比实际高 1.3--7.7 个数量级。可检出性由产物、架构与测量条件共同决定，时间证据不跨环境合并。% A-27 % A-28 % A-29

本作品不提出新检测原语，不主张新漏洞，不做密钥恢复；贡献在产物中心的证据组织与可审计判定。

\subsection{局限与未验证项}

本作品的结论受测量对象、测量环境、方法覆盖与证据完整度四类限制；这些限制决定结论能用到哪。

第一类是测量对象。M1 是 FIPS 203\cite{fips203} 初始公开草案（ipd）的实现，M2 是历史 Kyber768\cite{kyber2018} 快照，两者都不等同于最终标准；M5 才对应最终标准。M1 的计时数据是同一机器、同一锁频规程下的前期封存数据，不是与 M2--M5 同一批采集。M4 是 32 B 比较夹具，它的前后对照只说明检测链路能走通，不外推到完整解封装。这意味着，本作品的结论按对象分别成立，不能从一个对象推到“ML-KEM 实现”这个整体。

第二类是测量环境。全部时间结论来自一台 x86\_64 裸金属与 clang 18.1.3；本机未配置 isolcpus，SMT 开启，后台约 39 个服务在运行。最小步进 10 ns，32.311 ns 的效应只有 2--3 个分辨率单位。网络只有同机 veth 与 netem，没有真实远程链路；aarch64 只构成第二架构，不是第二台测量机，其容器未锁频、非裸金属。三台 WSL2 虚拟化复现只在机制层一致，时间数据没有进入任何 Holm 家族。这意味着，结论不能外推到其他 CPU、其他编译器、隔离更好的机器或真实远程。% A-30 % A-31 % A-32 % C-04 % C-05 % C-09

第三类是方法覆盖。M2 的主类对不是为除法操作数设计的，message-zero 与 message-ones 只让解密消息的位不同，没有控制除法的被除数；检测预算 25--67 ns 大于除法可能产生的效应。M2 在残留分支存在的情况下，预注册家族仍未拒绝，其原因没有分离。

工具一侧的覆盖同样有限。Binsec\cite{binsec} 是有界验证，只覆盖给定入口与范围；9×5 工具矩阵有 15 格未运行，Microwalk 只跑过 M3 一个对象；GCC 与 -O1/-Oz/-Og 三个级别未执行。这意味着，阴性的强度受方法覆盖限制；没有跑过的格子是未执行，不是未检出，方法覆盖之外不能推断。% C-01 % C-02 % C-10 % C-11 % C-15

第四类是证据完整度。真实缺陷的修复都采用上游补丁，自主修复只落在 M3、M4 两个人工夹具上；KyberSlash 官方修复版本的残留分支只定位、未修复、未复测。三对修复版本的输出摘要相同，只说明语义不变，不说明产物形态相同。题面要求中第 6、7 条按部分登记，缺口与原因见附录 A。面向 CI 的回归门禁未实现。这意味着，本作品证明的是机制消失与前后可复测；自主防护覆盖真实标准实现仍待补，部分项按部分引用。% C-07 % C-08 % A-35 % A-38

四类限制不互相抵消。一条结论同时受几类限制时，按最弱的一类引用；四类分开写，读者才能判断这条结论该信到哪一步。

\subsection{后续工作}

后续工作按“先补防护、再补功效、再补环境、收尾补工程化”排序；每条给出对应的局限、路径、成功判据与新代价。

第一项是 KyberSlash 官方修复版本的自主修复与复测，补上自主防护没有落到真实标准实现这一缺口。做法是把掩码选择移植到消息编码函数；判据是目标分支从反汇编消失、秘密标记报数由 256 降为 0、KAT 100/100 通过；计时复测视机时；代价是需要重新锁频采集。% C-07

第二项针对方法覆盖里主类对的缺口：M2 的输入类与样本量。做法是按除法操作数设计输入类，每类内部变化密文，并按 MDE 扩大样本；判据是能回答“残留分支存在却未检出”的原因；代价是机时增加，预注册家族需要重开。% C-15

第三项针对测量环境的量化与隔离缺口。做法是改用 rdtscp、配置 isolcpus、关闭 SMT；判据是单次分辨率与 32 ns 效应的比值改善；代价是需要独占机器。% C-04

第四项针对方法覆盖的编译器与级别缺口。做法是补 gcc 13.3.0 与 -O1/-Oz/-Og；判据是每一格都给出产物判定与计时判定，与现有 12 格同口径；代价是每格重新构建与采集。% C-01 % C-02

第五项针对测量环境的远程缺口。做法是从同机 veth 走到真实两端，测出本地与远程的检测预算比；判据是给出可复算的预算比与置信区间；代价是网络环境不可控，需要新的预注册。% C-05

第六项针对证据完整度的 CI 缺口。做法是把反汇编、Memcheck 与可变时延检查接入持续集成；判据是每次提交自动跑定位层并给出与基线的差异；代价是维护工具版本与基线，工具升级本身要重新登记。% C-08

六项的排序依据是依赖关系。第一项不补，自主防护的缺口仍然存在；第二项不补，M2 的阴性无法解释；环境与工程化放后，是因为它们不改变已有结论的强度，只扩大适用范围。
